\section{What an Event Costs}
\label{sec:ch14-what-an-event-costs}
\index{subscriptions!what a message costs|(}%

The plan in the last section says the router fetches a customer. It does not
say how often, and there are two plausible answers. The router could resolve
the author once, when the subscription opens, and reuse it; or it could resolve
it again for every message. The difference matters more than it sounds, because
a subscription is the one thing in a graph that can run for hours.

Counting is easy, because chapter~\ref{ch:strangling-the-monolith} moved the
request timeline into the shared library and all seven services now report one
line per request they serve. So: count the lines in every console,
open one subscription, count again, submit five reviews through the router,
count once more.

\begin{minted}[fontsize=\footnotesize,breaklines]{text}
                       catalog pricing inventory accounts reviews ordering nodes
opening the subscription     0       0         0        0       1        0     0
five events                  0       0         0        5       5        0     0
\end{minted}

The Accounts column is a step of both verification scripts, which count the
same log lines either side of their own run rather than trusting this table.
They open both of this chapter's subscriptions at once and write three reviews,
so the number they assert is six rather than five, and it is asserted exactly
rather than as a bound: a graph that started answering with fewer would be a
better graph and a wrong chapter.

\index{subscriptions!what opening one costs}%
The first row is the cheap surprise. Opening a subscription costs one request,
to one service, and touches nothing else in the graph. No entity fetch happens
at subscribe time, and it could not: there is no entity yet. The router has a
plan and a WebSocket and nothing to resolve.

\index{subscriptions!one entity fetch per message}%
The second row is the answer. Five messages, five requests to Accounts. The
five to Reviews are the five mutations rather than the events, because the
events come back down a connection that was already open. So the router
re-resolves the non-trigger half of the payload every single time, and
figure~\ref{fig:ch14-what-repeats} is the shape of that.

\begin{figure}[htbp]
  \centering
  % What happens once and what happens again. Referenced from chapter 14. Bare
% tikzpicture; the figure environment, caption and label live at the call site.
%
% The point of the picture is the asymmetry between the two rows. The trigger is
% a single request at the left edge and never recurs; every fetch under it is
% paid again for each message, so the cost of a subscription is set by how much
% of the payload the trigger's own subgraph cannot answer.
\begin{tikzpicture}[
    font=\small,
    dot/.style={circle, draw, fill=black, inner sep=1.1pt},
    hollow/.style={circle, draw, fill=white, inner sep=1.1pt},
    flow/.style={-{Stealth[length=2mm]}},
    note/.style={font=\scriptsize, gray, align=center, inner sep=1pt},
    lbl/.style={font=\scriptsize, align=right, inner sep=2pt}
  ]

  \draw[flow] (1.4,0) -- (10.4,0);
  \node[note, anchor=west] at (10.5,0) {time};

  % -- the trigger, once ------------------------------------------------------
  \node[lbl, anchor=east] at (1.3,0.9) {trigger};
  \node[dot] (trig) at (2.0,0.9) {};
  \node[note, anchor=south west, align=left] at (2.1,1.05)
    {one \code{subscribe} to \code{reviews},\\then nothing};

  % -- the events -------------------------------------------------------------
  \node[lbl, anchor=east] at (1.3,0) {events};
  \foreach \x in {3.4,4.8,6.2,7.6,9.0}{
    \node[hollow] at (\x,0) {};
  }

  % -- the fetch each one pays for --------------------------------------------
  \node[lbl, anchor=east] at (1.3,-1.1) {\code{accounts}};
  \foreach \x in {3.4,4.8,6.2,7.6,9.0}{
    \node[dot] (f\x) at (\x,-1.1) {};
    \draw[flow] (\x,-0.15) -- (\x,-0.95);
  }

  \node[note, anchor=north, align=center] at (6.2,-1.4)
    {one \code{\_entities} call per message, for the one field\\the subgraph
     holding the connection cannot answer};

  \draw[dashed, gray] (2.0,0.7) -- (2.0,-1.6);
  \node[note, anchor=north east, align=right] at (1.95,-1.6) {measured once};

\end{tikzpicture}

  \caption{The two rows are paid on different schedules. Opening the
  subscription is one request and never recurs; the entity fetch under it is
  paid again for every message that arrives. So the running cost of a
  subscription is decided not by how expensive it was to open but by how much
  of the payload the service holding the connection cannot answer on its own.}
  \label{fig:ch14-what-repeats}
\end{figure}

None of that makes subscriptions expensive, and reading it that way would be a
mistake. It is the ordinary cost of a field from another service, which
chapters~\ref{ch:enter-the-router} and~\ref{ch:entity-resolution-done-right}
both measured, attached to something that repeats. What is genuinely new is who
decides how often it is paid. A query is billed once to whoever sent it. A
subscription is billed per message, the client decides how many messages there
are by choosing what to subscribe to, and the schema decides what each one
costs. Those two decisions are made by different people, usually in different
teams, and neither of them sees the multiplication.

\subsection*{Which makes the selection set an operational decision}

\index{subscriptions!the selection set as a running cost}%
Drop \code{author} from the subscription document and the plan has no children
at all. Every message is then a single frame from Reviews and nothing else in
the graph is touched, forever. Add it back and the graph is doing an
\code{\_entities} call against Accounts on a schedule a client chose.

That is not an argument for narrow subscriptions. It is an argument for knowing
which fields on a subscription payload cross a boundary, because in a query
that knowledge buys you a paragraph in a design document and here it buys you a
throughput figure. A client that subscribes to a busy product and asks for one
field from a second service has quietly made that service's load proportional
to somebody else's write rate.

\index{subscriptions!what the subgraph is asked for}%
There is a version of this that is worse and Mosaic happens not to have. The
trigger returns one \code{Review} per message, so the entity fetch is one
representation. A subscription returning a list would batch, in the sense
chapter~\ref{ch:entity-resolution-done-right} measured, and the batch would be
rebuilt per message. Nothing here tested that, and the research note says so.

\subsection*{A word about the numbers I am not printing}

\index{measurement!what a heartbeat looks like}%
While measuring the above I nearly published a latency. The first bytes a
client receives from the router are usually a \code{:heartbeat} comment, and
across the runs recorded in the research note that first byte arrived anywhere
between one second and nearly five after the request. It is tempting to read
that as how long a subscription takes to become live. It is not. It is the
router's heartbeat schedule, and the subscription was live long before it.

The real gap is between a publish and the frame reaching the client, and this
chapter does not print it. Decision 62 in this book's spec says a timing gets a
committed script a reader can run twice rather than a number in a table, and
nothing here has one; the research note hands subscription timings to
chapter~\ref{ch:resilience-and-performance} with the rest of the load
questions. What is worth recording now is the mistake, because a plausible
number that measures the wrong thing is exactly the failure
chapter~\ref{ch:composition} caught three of.

\medskip
So the mechanism is understood and the cost is understood. What I had not
looked at yet was the connection itself, and that turned out to be where the
chapter's real finding was hiding.

\index{subscriptions!what a message costs|)}%
